% Generated by scripts/build_reports.py; edit reviewed Markdown inputs.
% Compile with XeLaTeX.
% Options for packages loaded elsewhere
\PassOptionsToPackage{unicode}{hyperref}
\PassOptionsToPackage{hyphens}{url}
\PassOptionsToPackage{dvipsnames,svgnames,x11names}{xcolor}
\documentclass[
  11pt,
]{article}
\usepackage{xcolor}
\usepackage[margin=1in]{geometry}
\usepackage{amsmath,amssymb}
\setcounter{secnumdepth}{-\maxdimen} % remove section numbering
\usepackage{iftex}
\ifPDFTeX
  \usepackage[T1]{fontenc}
  \usepackage[utf8]{inputenc}
  \usepackage{textcomp} % provide euro and other symbols
\else % if luatex or xetex
  \usepackage{unicode-math} % this also loads fontspec
  \defaultfontfeatures{Scale=MatchLowercase}
  \defaultfontfeatures[\rmfamily]{Ligatures=TeX,Scale=1}
\fi
\usepackage{lmodern}
\ifPDFTeX\else
  % xetex/luatex font selection
\fi
% Use upquote if available, for straight quotes in verbatim environments
\IfFileExists{upquote.sty}{\usepackage{upquote}}{}
\IfFileExists{microtype.sty}{% use microtype if available
  \usepackage[]{microtype}
  \UseMicrotypeSet[protrusion]{basicmath} % disable protrusion for tt fonts
}{}
\makeatletter
\@ifundefined{KOMAClassName}{% if non-KOMA class
  \IfFileExists{parskip.sty}{%
    \usepackage{parskip}
  }{% else
    \setlength{\parindent}{0pt}
    \setlength{\parskip}{6pt plus 2pt minus 1pt}}
}{% if KOMA class
  \KOMAoptions{parskip=half}}
\makeatother
\usepackage{longtable,booktabs,array}
\usepackage{caption}
\captionsetup[table]{skip=6pt}
\newcounter{none} % for unnumbered tables
\usepackage{calc} % for calculating minipage widths
% Correct order of tables after \paragraph or \subparagraph
\usepackage{etoolbox}
\makeatletter
\patchcmd\longtable{\par}{\if@noskipsec\mbox{}\fi\par}{}{}
\makeatother
% Allow footnotes in longtable head/foot
\IfFileExists{footnotehyper.sty}{\usepackage{footnotehyper}}{\usepackage{footnote}}
\makesavenoteenv{longtable}
\setlength{\emergencystretch}{3em} % prevent overfull lines
\providecommand{\tightlist}{%
  \setlength{\itemsep}{0pt}\setlength{\parskip}{0pt}}
\usepackage{amsmath,amssymb}
\usepackage{xurl}
\usepackage{fancyhdr}
\pagestyle{fancy}
\fancyhf{}
\fancyhead[L]{\small Pure Thought | Revised research specifications}
\fancyfoot[L]{\small 2026-09-17}
\fancyfoot[R]{\small\thepage}
\setlength{\headheight}{15pt}
\setlength{\emergencystretch}{3em}
\allowdisplaybreaks
\usepackage{bookmark}
\IfFileExists{xurl.sty}{\usepackage{xurl}}{} % add URL line breaks if available
\urlstyle{same}
% fallback for those not using the hyperref driver hyperxmp:
\makeatletter
\@ifundefined{xmpquote}{\newcommand{\xmpquote}[1]{#1}}{}
\makeatother
\hypersetup{
  pdftitle={11. Certified topological bands in an idealized Maxwell medium},
  pdfauthor={Pure Thought research portfolio},
  colorlinks=true,
  linkcolor={blue},
  filecolor={Maroon},
  citecolor={Blue},
  urlcolor={blue},
  pdfcreator={LaTeX via pandoc}}

\title{11. Certified topological bands in an idealized Maxwell medium}
\author{Pure Thought research portfolio}
\date{Revised 17 September 2026}

\begin{document}
\maketitle
\thispagestyle{empty}
\clearpage

\section{11. Certified topological bands in an idealized Maxwell medium}

\textbf{Revised specification --- 2026-09-11.}
\href{https://github.com/rick-stevens-ai/Pure-Thought/blob/main/reviews/11-Photonic-Topological-Crystals.md}{Detailed
critique} \ensuremath{\cdot}
\href{https://github.com/rick-stevens-ai/Pure-Thought/blob/main/archive/original-PRDs/11-Photonic-Topological-Crystals.md}{Original
description} \ensuremath{\cdot}
\href{https://github.com/rick-stevens-ai/Pure-Thought/blob/main/reviews/EVIDENCE-STANDARD.md}{Evidence
standard}

\subsection{Plain-language guide}

\subsubsection{The problem in everyday terms}

A repeating optical structure changes which light waves can travel
through it. With suitable nonreciprocal response, it may also support
light that travels along an interface in a preferred direction. This
problem asks whether a precisely defined ideal electromagnetic medium
really has the required band structure and topology. It separates a
mathematically guaranteed wave effect from the later task of finding a
material that realizes the assumed response.

\subsubsection{Key terms}

\begin{itemize}
\item
  \textbf{Maxwell equations} --- The equations governing classical
  electric and magnetic fields.
\item
  \textbf{Periodic medium} --- A medium whose properties repeat from
  cell to cell.
\item
  \textbf{Permittivity \ensuremath{\epsilon} and permeability
  \ensuremath{\mu}} --- Quantities describing how a medium responds to
  electric and magnetic fields; they can be direction-dependent tensors.
\item
  \textbf{Nonreciprocity} --- Response that breaks the usual interchange
  symmetry between source and receiver; here it provides the required
  time-reversal-breaking mechanism.
\item
  \textbf{Photonic band} --- An allowed family of electromagnetic wave
  frequencies across crystal momenta.
\item
  \textbf{Transverse constraint} --- The condition selecting physical
  wave fields consistent with the relevant divergence equation.
\item
  \textbf{Chern number} --- An integer characterizing the topology of an
  isolated group of wave bands.
\item
  \textbf{Interface mode} --- A wave concentrated near the boundary
  between two media.
\item
  \textbf{Lossless and nondispersive} --- Idealizations in which energy
  is not absorbed and the constitutive response is taken independent of
  frequency.
\end{itemize}

\subsubsection{Why this matters}

Controlling where and in which direction waves travel is a basic aim of
optical engineering. Topology offers one mechanism for robust interface
behavior under appropriate conditions. A certified Maxwell calculation
would also expose whether an apparent effect is caused by numerical
artifacts or by assumptions a real material cannot satisfy.

\subsubsection{What progress would mean}

The first meaningful achievement would be a validated ideal medium with
a proved gap and topological index. It could inform later waveguide
designs, but performance in a lossy, frequency-dependent device would
require a separate physical model.

\subsection{Core question}

Can a specified lossless periodic electromagnetic model support a
rigorously isolated band group with nonzero Chern number, and can its
interface modes be validated?

\subsection{Scope and assumptions}

Choose a two-dimensional periodic, lossless, frequency-independent
constitutive model with Hermitian positive \ensuremath{\epsilon}(r),
\ensuremath{\mu}(r), specified boundary conditions, and a mechanism
breaking physical time reversal. Treat the constitutive tensors as
mathematical inputs. Real-frequency dispersion, loss, and fabrication
feasibility are extensions requiring additional physical data.

Start by validating a reciprocal scalar dielectric benchmark with zero
total Chern number for a time-reversal-invariant band group; then
introduce an explicitly nonreciprocal tensor.

\subsection{Mathematical target}

In a formulation supporting the chosen model, solve

\[(\nabla +ik)\times \mu ^{-1}(r)(\nabla +ik)\times u=\omega ^{2}\epsilon (r)u,\]

in units with c=1, subject to the appropriate transverse constraint. Use
the \ensuremath{\epsilon}-weighted inner product or an equivalent
Hermitian formulation. Specify an isolated positive-frequency band
group; photons do not require an electron-like filling assumption. Prove
a frequency gap over the entire Brillouin zone and certify the projector
invariant in the correct metric.

\subsection{Required outputs and proof obligations}

Return ideal geometry and constitutive tensors, solver discretization,
eigenvalue enclosures, spurious-mode controls, gap and topology
certificates. Interface calculations must state both media and
termination. Report robustness only for an explicit perturbation class
and bound that preserves the required gap/symmetry.

\subsection{Validation and rejection controls}

Check homogeneous-medium dispersion, mesh/basis convergence,
Hermiticity, and transverse modes. Separate Galerkin error from k-space
interpolation error. Verify reciprocal controls cannot acquire a net
Chern number from numerical artifacts. Compare an interface spectrum
with the bulk-index difference without claiming immunity to arbitrary
disorder.

\subsection{Milestones and research extension}

Baseline: validated reciprocal eigenproblem. Strong: one certified ideal
nonreciprocal topological design. Research extension: a realistic
dispersive constitutive model or gap optimization within a bounded
geometry family. Fabrication drawings alone do not certify material
behavior.

Milestones are acceptance gates, not calendar promises. Before a run,
freeze the concrete instance/domain, tolerance, compute budget, and
permitted dependencies. A bounded search that fails to find a result
must preserve an unresolved outcome.

\subsection{Starting instruction}

Work on the baseline above. First make every input and convention
explicit, then validate the reference case. Build the checker alongside
the calculation. Label numerical evidence, finite-model certificates,
and continuum theorems separately; return the strongest justified result
and identify the exact unmet proof obligation.

\subsection{Primary references and their role}

\begin{itemize}
\tightlist
\item
  \href{https://journals.aps.org/prl/abstract/10.1103/PhysRevLett.100.013904}{Haldane
  and Raghu, Directional optical waveguides with broken time-reversal
  symmetry} --- Nonreciprocal media as the mechanism for photonic
  quantum-Hall-like edge modes.
\end{itemize}

References support definitions or limitations; the scope choices and
proposed milestones are editorial recommendations, not claims that these
results have already been obtained.

\subsection{Displayed formulation and claim boundary}

\href{https://github.com/rick-stevens-ai/Pure-Thought/blob/main/PRDs/11-Photonic-Topological-Crystals.md}{Canonical
specification}.

\[(\nabla+ik)\times\mu^{-1}(r)(\nabla+ik)\times u
=\omega^2\epsilon(r)u,\qquad
\langle u,v\rangle_\epsilon=\int_{\mathrm{cell}}u^*\cdot\epsilon v\,dr.\]

Use positive Hermitian constitutive tensors, the transverse constraint,
units \(c=1\), and an isolated positive-frequency band group. Certify
its gap and projector invariant in the weighted metric. Nonzero total
Chern number requires breaking the relevant physical time-reversal
symmetry; scalar reciprocal benchmarks provide rejection controls.

\subsection{Detailed review and rationale}

\href{https://github.com/rick-stevens-ai/Pure-Thought/blob/main/PRDs/11-Photonic-Topological-Crystals.md}{Revised
problem} \ensuremath{\cdot}
\href{https://github.com/rick-stevens-ai/Pure-Thought/blob/main/archive/original-PRDs/11-Photonic-Topological-Crystals.md}{Preserved
original} \ensuremath{\cdot}
\href{https://github.com/rick-stevens-ai/Pure-Thought/blob/f854d4e12b55c1f5b5c2a964b75f2d23496d981e/PRDs/11-Photonic-Topological-Crystals.md}{Pinned
upstream source}

Line references below refer to the preserved original at commit
\texttt{f854d4e12b55}, not the rewritten file. ``Blocking'' means the
original issue can invalidate the scientific target or claimed
certificate; ``major'' means it materially changes scope, inference, or
acceptance criteria. These are critiques of specifications, not results
of running the embedded research-code sketches.

\subsubsection{Detailed findings}

\begin{enumerate}
\def\labelenumi{\arabic{enumi}.}
\item
  \textbf{Blocking --- missing nonreciprocal constitutive physics
  (68--98).} A real scalar dielectric model does not implement the
  claimed gyromagnetic Chern phase. The relevant tensor and inner
  product must be included.
\item
  \textbf{Major --- electronic Z\ensuremath{{}_{2}} shortcut (95--98).}
  Ordinary photonic time reversal squares to +1; an electronic parity
  formula cannot be transplanted without an appropriate additional
  symmetry structure.
\item
  \textbf{Major --- ``exact'' Maxwell solution (21--25, 56--60).} Exact
  governing equations do not imply exact numerical eigenvalues.
  Discretization, spurious longitudinal modes, and continuum coverage
  need validation.
\item
  \textbf{Major --- fabrication and robustness (37--42, 740--754).}
  Generic dielectric rods and STL files do not establish nonreciprocity,
  telecom performance, or a universal 15\% disorder tolerance. The new
  problem certifies a mathematical medium first and identifies the extra
  evidence required for engineering claims.
\end{enumerate}

\subsubsection{What the revision preserves and changes}

The central motivation is retained, but the executable question is now:
\textbf{Can a specified lossless periodic electromagnetic model support
a rigorously isolated band group with nonzero Chern number, and can its
interface modes be validated?} The revised scope chooses a concrete
starting model; those choices are proposed research-design decisions,
not facts inferred from the original. The baseline can succeed without
resolving the full open research program.

The replacement also removes the original implementation sketches from
the active specification. They remain in the archive for comparison, but
have not been made executable or certified. New implementations should
derive their equations from the revised definitions rather than copy
unchecked prototypes.

\subsubsection{Acceptance criteria to use instead}

Baseline: validated reciprocal eigenproblem. Strong: one certified ideal
nonreciprocal topological design. Research extension: a realistic
dispersive constitutive model or gap optimization within a bounded
geometry family. Fabrication drawings alone do not certify material
behavior.

\subsubsection{Sources supporting the corrections}

\begin{itemize}
\tightlist
\item
  \href{https://journals.aps.org/prl/abstract/10.1103/PhysRevLett.100.013904}{Haldane
  and Raghu, Directional optical waveguides with broken time-reversal
  symmetry} --- Nonreciprocal media as the mechanism for photonic
  quantum-Hall-like edge modes.
\end{itemize}

\section{Evidence and execution standard}

This standard applies to the revised specifications. It is a proposed
contract for future research runs; this editorial revision does not
supply the scientific certificates themselves.

\subsection{What ``pure thought'' means here}

Use explicitly stated mathematical models and deductions without fitting
conclusions to experimental outcomes. Synthetic examples, exact model
inputs, analytic references, numerical search, and established software
are allowed. Physical constants, electronic-structure approximations,
constitutive laws, and empirical parameter sets must be identified when
used. A model-derived result is conditional on that model.

The original ban on legacy software conflicts with its proposed reliance
on NumPy, SciPy, CVXPY, nauty, and other libraries. Reimplementing every
optimizer and algebra system is not necessary for scientific
independence. Prefer a small independently checkable certificate and a
documented trusted base. A fresh implementation can be a project
objective, but must not be confused with a mathematical proof or novelty
claim.

\subsection{Four evidence levels}

{\def\LTcaptype{none} % do not increment counter
\begin{longtable}[]{@{}
  >{\raggedright\arraybackslash}p{(\linewidth - 4\tabcolsep) * \real{0.3333}}
  >{\raggedright\arraybackslash}p{(\linewidth - 4\tabcolsep) * \real{0.3333}}
  >{\raggedright\arraybackslash}p{(\linewidth - 4\tabcolsep) * \real{0.3333}}@{}}
\toprule\noalign{}
\begin{minipage}[b]{\linewidth}\raggedright
Label
\end{minipage} & \begin{minipage}[b]{\linewidth}\raggedright
What it establishes
\end{minipage} & \begin{minipage}[b]{\linewidth}\raggedright
What it does not establish
\end{minipage} \\
\midrule\noalign{}
\endhead
\bottomrule\noalign{}
\endlastfoot
Numerical evidence & A reproducible computation at stated precision,
sample size, and discretization & Exact truth, continuum coverage, or a
universal theorem \\
Finite-model certificate & An exact/validated statement about the
specified finite model or finite relaxation & An omitted tail,
thermodynamic limit, or physical realization \\
Conditional mathematical theorem & A proof over the stated domain with
every hypothesis and error controlled & Truth outside those assumptions
or a demonstrated physical system \\
Formal proof & A checked theorem in a named proof assistant, with
imported axioms/dependencies listed & Correct physical modeling merely
by virtue of formalization \\
\end{longtable}
}

Formalization is an optional extra deliverable unless the specific
project makes it mandatory. High precision, JSON output, solver status
and small residuals do not by themselves confer any of these proof
levels.

\subsection{Claim directions}

For minimization over a larger feasible set, the optimum is a lower
bound on the original minimum. A physical feasible point supplies an
upper bound. For maximization the directions reverse. Track feasibility
of the primal/dual witnesses separately from their objective values.

Infeasibility of an inner ansatz excludes only that ansatz.
Infeasibility of an outer relaxation can exclude the original problem if
the map into the relaxation is proved and all approximations preserve
necessity. Passing necessary tests means ``not excluded''; existence
needs a constructive sufficient argument.

A smallest observed example is not a minimality theorem. A complete
catalog must define the universe, equivalence relation, boundaries, and
proof that pruning misses nothing. A finite-size numerical trend does
not prove an asymptotic limit.

\subsection{Minimum run manifest}

Before calculation, record the problem ID and revision hash, exact
inputs and conventions, quantified domain, equivalence relation, target
observable, accepted evidence level, tolerances, and resource limits.
Record dependencies and precision/rounding mode. Where randomization is
used, preserve seeds, trial counts and raw sufficient statistics.
Resource limits are planning decisions, not evidence that a mathematical
solution cannot exist.

An output bundle should include:

\begin{itemize}
\tightlist
\item
  The statement being claimed, including all assumptions and
  quantifiers.
\item
  Construction or proof object, plus an independently runnable verifier
  and its trust assumptions.
\item
  Numerical/error ledger separating roundoff, discretization,
  truncation, optimization, sampling and modeling error where relevant.
\item
  Positive reference cases and negative controls that the verifier must
  reject.
\item
  Coverage: included domain, omitted domain, unresolved regions and
  termination reason.
\end{itemize}

These are requirements for future implementations, not a promise of
polynomial-size certificates for every problem.

\subsection{Checking specific certificate types}

\textbf{LP/SDP/SOS:} verify equality constraints, objective direction,
and cone feasibility using exact arithmetic or outward-rounded
intervals. Rationalizing a numerical solution can break PSD constraints.
Correct residuals with justified norm bounds; do not silently project
and reuse the old objective. Distinguish a proof of a polynomial
approximation from a proof for the original function.

\textbf{Interval numerics:} state function enclosures, domains and
outward rounding. A small root box needs existence and, if claimed,
uniqueness. Cover the complement for a global enumeration. Include
singularities and tails rather than discarding them numerically.

\textbf{SAT/SMT:} a solver's UNSAT label is not a certificate. Preserve
a proof supported by a checker and the proof that the encoding, symmetry
reduction and blocked solutions match the mathematical universe. DRAT
applies to an appropriate Boolean CNF proof, not automatically to every
real-arithmetic or SDP claim.

\textbf{Statistics:} specify the estimand, dependence assumptions,
confidence procedure and bias. Zero observed failures does not prove
zero failure probability. MCMC error bars based on estimated
autocorrelation are not generally rigorous mixing guarantees.

\subsection{Milestones and novelty}

Baseline means a working, independently validated result on a concrete
reference problem. Strong means a broader or sharper result with its
claimed evidence level. Research extension means a possible novel result
whose feasibility and significance need review. A negative theorem can
be valuable; an unsuccessful search can also be useful if its scope is
documented, but the two are not interchangeable.

Do not promise ``publication-quality'' in a fixed number of months.
Before claiming novelty, conduct a focused current literature comparison
against the exact theorem/model/metric. The sources in this revision
check key definitions and limitations; they are not an exhaustive
novelty audit across all 30 fields.

\end{document}
